% 2015 年重庆卷（理工农医类）第 5 题：某几何体的三视图（依固定扫描页重绘，只画拓扑与标号，不按比例）
\begin{tikzpicture}[>=stealth,scale=1.0]
  % 正视图：右侧 2×1 矩形，左侧接直角三角形（水平边 1、竖直边 1）
  \draw (0,0) rectangle (2,1);
  \draw (-1,0)--(0,1);
  \draw (-1,0)--(0,0);
  % 正视图尺寸标注
  \draw[<->] (-1,-0.4)--(0,-0.4);
  \node at (-0.5,-0.75) {$1$};
  \draw[<->] (0,-0.4)--(2,-0.4);
  \node at (1,-0.75) {$2$};
  \draw[<->] (2.28,0)--(2.28,1);
  \node at (2.62,0.5) {$1$};
  \node at (0.5,-1.25) {正视图};
  % 左视图：半径 1 的半圆，竖直半径与两条弦
  \draw (4.1,0)--(6.1,0);
  \draw (4.1,0) arc (180:0:1);
  \draw (5.1,0)--(5.1,1);
  \draw (5.1,1)--(4.1,0);
  \draw (5.1,1)--(6.1,0);
  \node at (5.1,-1.0) {左视图};
  % 俯视图：2×2 矩形，左侧接等腰三角形（底 2、高 1）
  \draw (0,-4) rectangle (2,-2);
  \draw (-1,-3)--(0,-2)--(0,-4);
  \draw (-1,-3)--(0,-3);
  % 俯视图尺寸标注（深度方向 1+1）
  \draw[<->] (0.3,-2)--(0.3,-3);
  \node at (0.62,-2.5) {$1$};
  \draw[<->] (0.3,-3)--(0.3,-4);
  \node at (0.62,-3.5) {$1$};
  \draw[<->] (0,-4.35)--(2,-4.35);
  \node at (1,-4.75) {$2$};
  \node at (0.5,-5.35) {俯视图};
\end{tikzpicture}
